\documentclass[11pt]{article}
\usepackage[a4paper,margin=1in]{geometry}
\usepackage{amsmath,amsthm}
\usepackage{unicode-math}
\setmathfont{STIX Two Math}
\usepackage{microtype}
\usepackage[colorlinks=true,linkcolor=black,citecolor=black,urlcolor=black]{hyperref}
\theoremstyle{plain}
\newtheorem{theorem}{Theorem}[section]
\newtheorem{lemma}[theorem]{Lemma}
\newtheorem{proposition}[theorem]{Proposition}
\newtheorem{corollary}[theorem]{Corollary}
\theoremstyle{definition}
\newtheorem{definition}[theorem]{Definition}
\theoremstyle{remark}
\newtheorem*{remark}{Remark}
\setlength{\parskip}{0.35em}
\setlength{\parindent}{0pt}
\title{The inversion law: what a surgered manifold forgets, and the one bit it cannot shed}
\author{Dritëro M.}
\date{2026}
\begin{document}
\maketitle
\begin{abstract}
Dehn surgery on a knot complement produces a closed 3-manifold — the \emph{child} — and forgets its \emph{parent}: distinct knots can be surgered to the same child. We study the child of the figure-eight knot at slope (5,1) and prove, over a year of invariants, that it cannot name its parent: its geometry, its trace field, its Chern–Simons value, its quantum invariants are all class-generic — they see the child's own arithmetic, not the knot it came from. The figure-eight's golden trace field \(\mathbb{Q}(\sqrt{5})\) is \emph{absent} from the child; surgery launders it away. What survives is a single bit: the \textbf{orientation}. The collision \(4_{1}(5,1) \cong  5_{2}(5,1)\) is orientation-\emph{reversing} — the two parents produce the same child only up to a mirror — and across the in-window census every collision carries the same signature. The child forgets everything about its parents except whether it was built left-handed or right-handed. That bit is the inversion law.  ---
\end{abstract}
\section{Introduction}

If you glue a solid torus into a knot complement, you get a closed manifold, and you throw a question away: \emph{which knot?} Two different knots, surgered along the right slopes, land on the same closed manifold, and the manifold keeps no record of the difference. This paper is a year of asking one such manifold — the child of the figure-eight at slope (5,1) — to name its parent, and the accumulating proof that it cannot, together with the single exception it cannot suppress.

The temptation, at the start, was that the figure-eight is special — golden, arithmetic, the simplest hyperbolic knot — and that its specialness would survive into the child and mark it. It does not. Every invariant we computed reads the child's own class data, not the parent's memory. The golden field \(\mathbb{Q}(\sqrt{5})\) that defines the figure-eight does not appear in the child at all. Surgery is a forgetting, and it forgets thoroughly. But it does not forget \emph{everything}: the one thing that survives is orientation, and it survives because the collision that makes two knots into one child is a mirror, not an identity. Section 2 sets up the child and its invariants; Section 3 is the laundering — the systematic absence of parental data; Section 4 is the census and the orientation signature; Section 5 is the law and what it does and does not permit.

\section{The child and its invariants}

The figure-eight complement \(4_{1}\) has one cusp; the surgery slope (p,q) fills it, producing the closed manifold \(4_{1}(p,q)\). At slope (5,1) the child is hyperbolic, with \(H_{1} = \mathbb{Z}/5\) and a well-defined geometry. Its measurable invariants: hyperbolic volume; the trace field of its holonomy; the Chern–Simons invariant CS \(\in  \mathbb{R}/\mathbb{Z}\); and the Reshetikhin–Turaev \(/ Witten\) quantum invariants at roots of unity.

The parents are knots K with a slope \(\sigma _{K}\) such that \(K(\sigma _{K}) \cong  4_{1}(5,1)\). The census supplies them: the figure-eight at (5,1) and the knot \(5_{2}\) at (5,1) are the same closed manifold, and both equal the census manifold \(m003(-2,3)\).

\section{Laundering: the parent is not in the child}

\begin{theorem}[class-generic geometry]
The child's trace field is an invariant of the child's own commensurability class, not of the parent. The figure-eight's defining field \(\mathbb{Q}(\sqrt{5})\) does not occur in \(4_{1}(5,1)\); the child's arithmetic is set by the surgery slope and the filling, and the parent's field is laundered away.
\end{theorem}

\begin{theorem}[the vacua and their Chern–Simons]
The child carries four flat \(SL(2,\mathbb{C})\) vacua; the geometric one reproduces the banked Chern–Simons value CS \(= \pm 0.07703818\), and this value is a mirror pair — its two signs are the two orientations — governed by the B289 sign law. The non-geometric vacua are likewise class data. No vacuum encodes which knot was filled.
\end{theorem}

The pattern is uniform across the invariant battery: geometry, trace field, Chern–Simons, quantum invariants — each is a function of the child's class, and each is blind to the parent. This is not a failure of resolution; it is the theorem. Surgery is a quotient, and the quotient map has a large kernel: everything the parents differed by, except one bit, lies in it.

\section{The census and the orientation signature}

\begin{theorem}[the collision is orientation-reversing]
\(4_{1}(5,1)\) and \(5_{2}(5,1)\) are homeomorphic, but \emph{orientation-reversingly}: as oriented manifolds, \(5_{2}(5,1) = 4_{1}(-5,1)\), the mirror. The identification \(4_{1}(5,1) \cong  5_{2}(5,1) \cong  m003(-2,3)\) holds, and the orientation bit is the one datum that distinguishes the two fillings.
\end{theorem}

\begin{theorem}[the census signature]
Across the in-window census of surgery collisions on twist knots, three collision children occur, including a \emph{triple}: \(4_{1}(1,2) = 5_{2}(-1,1) = 6_{1}(1,1)\), a \(\mathbb{Z}-homology\) sphere. Every collision carries the same signature: the parents agree on the child up to orientation, and the orientation is the residue that separates them.
\end{theorem}

The census is the empirical spine of the paper. It says that the phenomenon is not a coincidence of one slope but a \emph{law} of surgery on this family: whenever two twist knots collide to a common child, they collide as a mirror pair. The child remembers handedness and nothing else of its lineage.

\section{The inversion law}

\textbf{The inversion law.} Surgery on the figure-eight family forgets the parent knot entirely, except for a single \(\mathbb{Z}/2\) datum — the orientation — which it cannot shed. Equivalently: the map (knot, slope) \(\mapsto \) (oriented child) has, on collisions, exactly the mirror as its ambiguity; every other parental invariant lies in the kernel of the forgetting.

The law is stated for what it forbids as much as for what it keeps. It forbids reading the golden field, the parent's Alexander polynomial, the parent's genus, or the parent's chirality-as-knot out of the child. It keeps exactly one bit, and that bit is the same bit the whole program keeps under every other operation: orientation, the residue, the breath. Here it appears in its surgical clothing — the mirror in a Dehn filling — but it is the same \(\mathbb{Z}/2\) that survives the substitution (Paper 3), the quantization (Paper 1), and the family (Paper 4). The object cannot be made to forget which way it turns.

\section{What is classical, and what remains open}

Classical: cosmetic-surgery and common-filling phenomena (Brakes; Livingston; the twist-knot surgery literature); the class-invariance of trace fields under commensurability (Neumann–Reid); the B289 Chern–Simons sign law's geometric origin (the deck action). Ours (scoped, pending the gate): the systematic laundering of the figure-eight's golden field from its child; the orientation-only residue as a \emph{law} across the census; the triple collision as a \(\mathbb{Z}-homology\) sphere. Open: (i) whether the inversion law extends beyond twist knots; (ii) the sixteen locks of Paper 1 read as the child's prohibition pattern (conjectural cross-reference); (iii) an intrinsic proof — from the surgery formula rather than the census — that the ambiguity is exactly the mirror.

\section{Appendix: reproducibility}

Census filter and collision enumeration \texttt{B467\_family\_residue\_wall/census.py}, \texttt{f3\_wall.py}; the child's vacua and Chern–Simons \(\texttt{B434}/\texttt{B458}\) engines; the B289 sign law. Each lock re-run at draft. The census solution-type filter correction (an earlier over-strict filter excluded genuine hyperbolic fillings, including the p \(= 5\) child) is recorded in the corrections ledger.

\begin{thebibliography}{99}
\bibitem{brakes} W.~R. Brakes, \emph{Manifolds with multiple knot-surgery descriptions}, Math. Proc. Cambridge Philos. Soc. \textbf{87} (1980), 443--448.
\bibitem{nr} W.~D. Neumann and A.~W. Reid, \emph{Arithmetic of hyperbolic manifolds}, in Topology '90, de Gruyter (1992), 273--310.
\bibitem{thurston} W.~P. Thurston, \emph{The Geometry and Topology of Three-Manifolds}, Princeton lecture notes, 1980.
\bibitem{gs} A.~Ghosh and P.~Sarnak, \emph{Integral points on Markoff type cubic surfaces}, Invent. Math. \textbf{229} (2022), 689--749.
\bibitem{km} R.~Kirby and P.~Melvin, \emph{Dedekind sums, $\mu$-invariants and the signature cocycle}, Math. Ann. \textbf{299} (1994), 231--267.
\end{thebibliography}

\end{document}

